\documentclass[10pt,twocolumn,letterpaper]{article}
\usepackage{cvpr}
%% This file contains a number of tweaks that are typically applied to the main document.
%% They are not enabled by default, but can be enabled by uncommenting the relevant lines.


%%
%% Official support for 2-column teaser figure
%%
\usepackage{cuted} %< for strip environment

%%
%% For auto-referencing of figures (see fig/teaser.tex and \label{} therein) 
%%
\usepackage{currfile} %< for \currfilebase macro

%%
%% caption and (global) spacing overrides
%%
\usepackage{caption} %< for "\captionof" commands (teaser)
% \setlength{\abovecaptionskip}{.5em} % space above caption
% \setlength{\belowcaptionskip}{1em} % space below caption

%%
%% Inline annotations; for predefined colors, refer to "dvipsnames" in the xcolor package:
%% https://tinyurl.com/overleaf-colors
%%
\newcommand{\red}[1]{{\color{red}#1}}
\newcommand{\todo}[1]{{\color{red}#1}}
\newcommand{\TODO}[1]{\textbf{\color{red}[TODO: #1]}}
%%
%% Disable for camera ready / submission by uncommenting these lines:
%%
% \renewcommand{\TODO}[1]{}
% \renewcommand{\todo}[1]{#1}

%%
%% Work harder in optimizing text layout. Typically shrinks text by 1/6 of page, enable
%% it at the very end of the writing process, when you are just above the page limit.
%%
% \usepackage{microtype}

%%
%% Fine-tune paragraph spacing.
%%
% \renewcommand{\paragraph}[1]{\vspace{.5em}\noindent\textbf{#1.}}

%%
%% Globally adjusts space between figure and caption.
%%
% \setlength{\abovecaptionskip}{.5em}

%%
%% Allows "the use of \paper to refer to the project name"
%% with automatic management of space at the end of the word.
%%
% \usepackage{xspace}
% \newcommand{\paper}{ProjectName\xspace}

%%
%% Commonly used math definitions.
%%
% \DeclareMathOperator*{\argmin}{arg\,min}
% \DeclareMathOperator*{\argmax}{arg\,max}

%%
%% Tighten underline.
%%
% \usepackage{soul}
% \setuldepth{foobar}

%%
%% Reconfigure paragraph (tweaks spacing)
%%

\definecolor{cvprblue}{rgb}{0.21,0.49,0.74}
\usepackage[pagebackref,breaklinks,colorlinks,allcolors=cvprblue]{hyperref}
\def\paperID{}
\def\confName{CVPR}
\def\confYear{2026}
\input{generated/assignment2/selected_response_metrics.tex}
\input{generated/ce7454/status.tex}
\title{Direction Alignment for One-Step Text-to-Image Distillation}
\author{Le Duc Dong \quad Nguyen Huu Thuat \quad Pham Viet Hoang \quad La Nguyen Gia Hy\\CE7454 Project 2}
\begin{document}
\maketitle
\begin{abstract}
Compressing a diffusion teacher into one denoiser call is difficult under a
small training budget. We investigate a direction-alignment extension to
distribution matching distillation (DMD). For two noises sharing a caption,
the added loss aligns the student latent displacement with the corresponding
teacher displacement while retaining DMD's perceptual regression. A rank-32
Stable Diffusion 1.5 student trained independently for 36,000 updates obtains
FID \ASelectedFID{}, IS \ASelectedIS{} and CLIP \ASelectedCLIP{} on all 5,000
COCO-2017 validation captions at one network function evaluation. The original
DMD+LPIPS reference has FID 29.94, IS 34.12 and CLIP 31.10. The added
constraint aligns the direction of paired latent changes without
requiring their magnitudes to match. Different teacher banks and initialization histories limit
attribution of the main comparison to the added loss alone.
\end{abstract}

\section{Problem and motivation}
Diffusion models refine an image through repeated denoising, making good
generation expensive. Latent diffusion reduces each step's cost by operating
in a VAE representation~\cite{rombach2022ldm}; distillation reduces the number
of steps. Our project asks whether a one-step generator with limited LoRA
capacity can better preserve the teacher's response to changes in noise.

DMD combines diffusion-score matching and regression to teacher-generated
images~\cite{yin2024dmd}. Its LPIPS term anchors individual outputs but does
not explicitly align the change between two outputs. We add a direction-only
constraint on that change. This reuses the paired regression examples and
adds no student parameters or inference calls. The difficulty is balancing
the extra training signal against image fidelity and caption alignment.

\begin{figure*}[t]
\centering
\includegraphics[width=\textwidth]{generated/assignment2/main_samples.pdf}
\caption{Native samples with identical row prompts and initial noise. All
students have 36,000 updates and use one step / one NFE. The teacher uses
64 DDIM steps / 128 NFE. Ours is the independently trained model evaluated
in Table~\ref{tab:quality}. Dog and lighthouse prompts are preset; teddy
bears are a selected validation example. These examples do not determine
the full-validation metric ranking.}
\label{fig:samples}
\end{figure*}

\section{Related work}
Progressive distillation compresses trajectories through two-to-one
transitions~\cite{salimans2022progressive}; consistency distillation learns
a shared endpoint along a trajectory~\cite{song2023consistency}. SwiftBrush
uses variational score distillation without a paired-image bank, while
InstaFlow simplifies noise--image couplings through
reflow~\cite{nguyen2024swiftbrush,liu2024instaflow}.

Perceptual regression is part of original DMD, and response alignment has
prior art. GAD matches local teacher and student responses with vector
error~\cite{huang2026gad}. We test angular alignment between independent
cached pairs in a restricted SD1.5/LoRA setting. The contribution is the
specific direction-only objective and its experimental study, rather than
the invention of perceptual regression, response matching or cosine loss.

\section{Method}
\paragraph{One-step DMD backbone.}
The student $G_\theta(z,c)$ maps Gaussian noise and a caption to a clean
latent with one conditional U-Net call. LoRA adapts its attention projections
at rank 32~\cite{hu2022lora}. A frozen teacher and a learned fake scorer
predict clean latents $\hat x_0^T$ and $\hat x_0^F$ from the same noised
student output $x$. The normalized distribution-matching direction is
\begin{equation}
g=\frac{\hat x_0^F-\hat x_0^T}
{\max\{\operatorname{mean}(|x-\hat x_0^T|),\epsilon\}}.
\end{equation}
A stopped-target surrogate injects $g$ into the generator. A separate
epsilon-denoising update fits the fake scorer on detached generated samples.
The VAE and text encoder remain frozen, and teacher guidance is baked into
the student; inference uses a single conditional branch.

\paragraph{Direction alignment.}
Two cached noises $z_1,z_2$ share caption $c$ and teacher endpoints
$T(z_1,c),T(z_2,c)$. Define
\begin{align}
d_S&=G_\theta(z_2,c)-G_\theta(z_1,c),\nonumber\\
d_T&=T(z_2,c)-T(z_1,c),\nonumber\\
L_{\rm angle}&=1-\frac{\langle d_S,d_T\rangle}{\|d_S\|_2\|d_T\|_2},\\
L&=L_{\rm DMD}+0.25L_{\rm LPIPS}+\lambda L_{\rm angle}.
\end{align}
Numerical guards handle vanishing norms. $L_{\rm LPIPS}$ averages LPIPS-VGG
between decoded, clamped student and teacher RGB images~\cite{zhang2018lpips}.
The teacher targets are detached; both student branches receive gradients.
Pairs with teacher displacement norm at most $10^{-6}$ are excluded, and
the student norm is clamped below at $10^{-6}$ in the denominator. The
displacements use flattened FP32 latents.

The paired DMD outputs already provide both student endpoints, so the added
term requires no extra teacher or generator forward pass relative to that
paired update. We fix $\lambda=0.298494$ from the first update of the
independent run. This study does not establish an optimal angular weight
within the 36,000-update budget.

\begin{figure*}[t]
\centering
\includegraphics[width=.96\textwidth]{generated/assignment2/response_alignment_method.pdf}
\caption{Direction-alignment design. Cached teacher endpoints share a caption
and two independent noises. Their latent displacement guides the direction
of the student displacement; both student branches receive gradients.
DMD and LPIPS remain in the objective. The vectors are schematic.}
\label{fig:method}
\end{figure*}

\section{Experimental protocol}
\paragraph{Data and metrics.}
We use COCO-2017 and select the lowest-ID caption per
image~\cite{chen2015coco}. Two disjoint pools of 50,000 and 68,287 captions
cover the training split. All 5,000 validation images are held out from
student training. Each final evaluation generates one $512\times512$ image
per validation caption using the same noise stream, seed 12345, and the
same resizing and center-cropping protocol. We report FID, Inception Score
and CLIP ViT-B/32 similarity, defined as mean $100\max(\cos(I,c),0)$.

\paragraph{Training and comparison.}
Ours starts with pretrained SD1.5 and fresh rank-32 student/scorer adapters,
receiving 18,000 updates per caption pool. Effective batch is 16 via
microbatch two and eight accumulation passes. AdamW uses learning rate
$10^{-4}$, 100-step warmup, cosine decay, weight decay 0.01, gradient clip 1
and EMA 0.995. The scorer receives one update per generator update. Teacher
CFG is 7.5; student inference CFG is 1. Each phase has 5,000 cached teacher
outputs from 2,500 captions and two noises each: 10,000 outputs and
1,280,000 teacher NFEs in total.

All standard baselines use rank-32 generators, 36,000 updates and one-NFE
inference. The local consistency model uses a VP boundary blend and latent
L2; InstaFlow uses VP epsilon-to-velocity conversion and cached couplings.
SwiftBrush retains teacher CFG 4.5 and FP32; other recipes use CFG 7.5 and
BF16. Original DMD+LPIPS uses a 5,000-output bank and an earlier training
history. Equal generator updates therefore do not equalize supervision or
total compute. \CEMatchedControlStatus{}

\section{Results and analysis}
\begin{table}[t]
\centering\small
\input{generated/ce7454/quality.tex}
\caption{Full COCO-2017 validation: 5,000 generated and 5,000 real images.
All students use one NFE; the teacher uses 128. Baseline names denote our
local adapter recipes. Bold marks the best student metric.}
\label{tab:quality}
\end{table}
\input{generated/ce7454/findings.tex}

\paragraph{What the angular term changes.}
Let $u=d_S/\|d_S\|$ and $v=d_T/\|d_T\|$. Above the norm guard,
\begin{equation}
\nabla_{d_S}L_{\rm angle}
=-\frac{v-(u^\top v)u}{\|d_S\|}.
\end{equation}
The gradient removes the component parallel to $d_S$ and rotates the
student response toward the teacher response. The two endpoint gradients
have opposite signs, so their sum is zero in latent coordinates. Their
magnitude is proportional to $|\sin\phi|/\|d_S\|$, where $\phi$ is the angle
between the displacements. Small student displacements therefore require
the numerical guard. These properties concern this loss alone; shared
network parameters and Adam scaling can also change magnitude and midpoint.

For $r=\|d_S\|/\|d_T\|$, a full relative vector error decomposes as
\begin{equation}
\frac{\|d_S-d_T\|^2}{\|d_T\|^2}
=(r-1)^2+2r(1-u^\top v).
\end{equation}
Our angular term separates orientation from the radial penalty and its
$r$-dependent weighting. This leaves output scale and appearance to the
existing DMD and LPIPS objectives. The choice is useful when relative
response direction is the desired constraint, but it does not force the
correct amount of variation. Independent noise pairs define finite changes;
they do not estimate a local Jacobian or establish improved diversity.

\paragraph{Quality and text control.}
Lower FID accompanies lower IS and CLIP, so the result is a trade-off in
measured quality rather than a uniform improvement. FID compares marginal
image-feature statistics and can improve without better caption binding.
The angular objective constrains noise response under a fixed caption;
it does not explicitly supervise object counts or relations. These
distinctions explain why the text-control figures are useful alongside
the aggregate metrics.

\begin{figure}[t]
\centering
\includegraphics[width=\linewidth]{generated/assignment2/prompt_control.pdf}
\caption{Ours under fixed noise: four animals crossed with reading, drinking
and eating prompts. All twelve combinations are shown at one NFE.}
\label{fig:control}
\end{figure}
\begin{figure}[t]
\centering
\includegraphics[width=\linewidth]{generated/assignment2/prompt_interpolation.pdf}
\caption{Ours with dog-to-cat hidden-state interpolation at weights
$0,0.25,0.5,0.75,1$. Each row fixes its initial noise. Images are native
generator outputs rather than pixel blends.}
\label{fig:interpolation}
\end{figure}
In Figure~\ref{fig:control}, animal identity and activity change while
the cafe setting remains recognizable under fixed noise.
Figure~\ref{fig:interpolation} varies text hidden states, a different input
from the noise pairs used in training. It demonstrates conditioning
behavior of the same one-step model, rather than directly testing the
angular constraint. All images are native outputs; the selected examples
complement the full-set metrics.

\section{Limitations and conclusion}
Resource constraints limit each main student to 36,000 updates and one
training seed. Published budgets differ substantially: SwiftBrush reports
65,000 updates, and consistency-model experiments include
600,000--1,000,000 iterations~\cite{nguyen2024swiftbrush,song2023consistency}.
Architecture, batch size and teacher cost prevent direct compute equivalence.
The small paired bank, validation reuse and unknown teacher-pretraining
overlap with COCO limit generalization claims.

The project shows that direction-aligned DMD can produce useful one-step
images under restricted adapter capacity. Its independently trained model
reaches FID \ASelectedFID{} with lower IS and CLIP than the original DMD
reference, without adding inference calls or trainable parameters.
\CEControlConclusion{} Dedicated latency measurements
would also need isolated GPU workloads.

{\small
\bibliographystyle{ieeenat_fullname}
\bibliography{assignment2_references}
}
\end{document}
